\chapter{Waring's Problem}\label{ch:waring}
\begin{center}\itshape
Every number is a sum of four squares, of nine cubes,\\
of nineteen fourth powers, ``and so on''.\\
How many ways are there, once the number is large?
\end{center}

Edward Waring, Lucasian professor at Cambridge, stated in his \emph{Meditationes Algebraicae} of 1770 that every natural number is a sum of four squares, of nine cubes, of nineteen fourth powers, ``and so on''. Lagrange proved the four squares the same year. The cubes waited until 1909, when Wieferich showed that nine suffice, and in the same year Hilbert proved that for every $k$ some finite number $g(k)$ of $k$-th powers suffices. Hilbert's proof gave no reasonable bound.

Here is a small count to hold on to. The number $79$ is $4\cdot16+15\cdot1$, and since $81>79$ the only fourth powers available are $1$ and $16$. Four sixteens and fifteen ones use nineteen fourth powers, and no choice uses fewer. So $g(4)\ge19$, and Waring's ``nineteen'' is sharp. Small numbers are stubborn for silly reasons. The question that has a beautiful answer is about large numbers: let $G(k)$ be the least $s$ such that every \emph{sufficiently large} integer is a sum of $s$ positive $k$-th powers.

In 1920 Hardy and Littlewood, in the first of their papers \emph{Partitio Numerorum}, wrote the number of representations of $n$ as a sum of $s$ $k$-th powers as a Fourier coefficient of the $s$-th power of a Weyl sum. The circle splits into \emph{major arcs}, near fractions with small denominators, where the Weyl sum can be computed, and \emph{minor arcs}, where Weyl's inequality makes it small. Volume~5, Chapter~12, used the same picture for partitions. In 1938 Hua Loo-Keng, who had taught himself mathematics in his father's shop in Jintan and had come to Cambridge in 1936 to work with Hardy, found a mean value estimate that makes the asymptotics work for $2^k+1$ powers.

This chapter proves Weyl's inequality for degree $k$ and Hua's lemma, both by Weyl's differencing, then the asymptotic formula for every $s\ge 2^k+1$. The main term contains a \emph{singular series}, a product of local densities, one for each prime. We prove that it is bounded below by a positive constant when $s\ge4k$ as well, so that $G(k)\le2^k+1$ for $k\ge4$. There are no primes in this chapter. It is the circle method in its cleanest form, and Chapter~\ref{ch:threeprimes} runs the same machine on the primes.

\begin{reminder}[from Volumes 1, 5, 6 and 8, and Chapter~\ref{ch:expsums}]
\emph{Geometric sums} (Volume~6, Lemma~11.1). For an interval $I$, $\bigl|\sum_{n\in I}e(\beta n)\bigr|\le\min\bigl(|I|+1,\frac1{2\|\beta\|}\bigr)$, where $e(t)=e^{2\pi it}$ and $\|t\|$ is the distance from $t$ to the nearest integer.
\emph{Sums of minima} (Volume~6, Lemma~11.6). If $\bigl|\alpha-\frac aq\bigr|\le\frac1{q^2}$ with $(a,q)=1$, then $\sum_{1\le x\le X}\min\bigl(N,\frac1{\|\alpha x\|}\bigr)\le\bigl(\frac Xq+1\bigr)(4N+2q+2q\log q)$ for real $X,N\ge1$. For every $\varepsilon>0$, the number of divisors satisfies $\tau(m)\ll_\varepsilon m^\varepsilon$.
\emph{Dirichlet's approximation} (Lemma~\ref{lem:dirichlet}). For real $\alpha$ and $Q\ge1$ there are coprime $a,q$ with $1\le q\le Q$ and $\bigl|\alpha-\frac aq\bigr|\le\frac1{qQ}$.
\emph{The Beta integral} (Volume~5, Theorem~5.11). $\int_0^1 t^{x-1}(1-t)^{y-1}\dd t=\frac{\Gamma(x)\Gamma(y)}{\Gamma(x+y)}$ for $x,y>0$.
\emph{Fourier inversion} (Volume~8, Theorem~10.2). If $h\in L^1(\R)$ and $\hat h\in L^1(\R)$, then $h(x)=\int\hat h(\xi)e(x\xi)\dd\xi$ for almost every $x$. The right side is continuous in $x$, so if $h$ is continuous the equality holds at every $x$.
\end{reminder}

\section{Weyl's inequality and Hua's lemma}

For Waring's problem the generating function is the Weyl sum
\[
f(\alpha)=\sum_{1\le x\le P}e(\alpha x^k),\qquad k\ge2,\ P\ge1,
\]
and the tool is differencing, the move behind Weyl's equidistribution theorem (Volume~6, Theorem~11.4), now made quantitative. For a function $\psi$ on the integers and $h\in\Z$, let $\Delta_h\psi(x)=\psi(x+h)-\psi(x)$, and for $\mathbf h=(h_1,\dots,h_j)$ let $\Delta_{\mathbf h}=\Delta_{h_1}\cdots\Delta_{h_j}$. Expanding, $\Delta_{\mathbf h}\psi(x)=\sum_{S\subseteq\{1,\dots,j\}}(-1)^{j-|S|}\psi(x+h_S)$ with $h_S=\sum_{i\in S}h_i$. For an interval $I\subseteq[1,P]$ let $T(\psi,I)=\sum_{x\in I\cap\Z}e(\psi(x))$, let $H=\Z\cap\left]-P,P\right[$, and let $I_{\mathbf h}$ be the set of $x$ with $x+h_S\in I$ for all $S$, an interval contained in $I$.

\begin{lemma}[Weyl differencing]\label{lem:weyldiff}
For $j\ge1$,
\[
|T(\psi,I)|^{2^j}\le(2P)^{2^j-j-1}\sum_{\mathbf h\in H^j}T(\Delta_{\mathbf h}\psi,I_{\mathbf h}),
\]
and the right side is real and non-negative.
\end{lemma}
\begin{proof}
For any interval $J\subseteq[1,P]$ and function $\phi$, writing $x=y+h$,
\[
|T(\phi,J)|^2=\sum_{x,y\in J}e(\phi(x)-\phi(y))=\sum_{h\in H}T(\Delta_h\phi,J\cap(J-h)).\tag{$\ast$}
\]
For $j=1$ this is the claim. Suppose the claim holds for $j$. By $(\ast)$ with $\phi=\Delta_{\mathbf h}\psi$ and $J=I_{\mathbf h}$, summed over $\mathbf h\in H^j$, the right side for $j+1$ is $(2P)^{2^{j+1}-j-2}\sum_{\mathbf h\in H^j}|T(\Delta_{\mathbf h}\psi,I_{\mathbf h})|^2\ge0$; and the right side for $j$ is $\sum_{\mathbf h\in H^{j-1}}|\cdots|^2\ge0$ by the same identity (for $j=1$ it is $|T(\psi,I)|^2$). Squaring the inequality for $j$ and applying Cauchy--Schwarz over the at most $(2P)^j$ vectors $\mathbf h$,
\[
|T(\psi,I)|^{2^{j+1}}\le(2P)^{2^{j+1}-2j-2}(2P)^j\sum_{\mathbf h\in H^j}|T(\Delta_{\mathbf h}\psi,I_{\mathbf h})|^2,
\]
which is the claim for $j+1$.
\end{proof}

For $\psi(x)=\alpha x^k$ the differences have a simple shape. For $j\le k$ and $x\in I_{\mathbf h}$,
\[
\Delta_{\mathbf h}x^k=h_1\cdots h_j\,p_j(x;\mathbf h),\qquad p_j(x;\mathbf h)=\int_{[0,1]^j}\frac{k!}{(k-j)!}(x+\mathbf t\cdot\mathbf h)^{k-j}\dd\mathbf t,\tag{$\dagger$}
\]
by the fundamental theorem of calculus applied in each variable $t_i$ in turn. The point $x+\mathbf t\cdot\mathbf h$ is multi-affine in $\mathbf t$, so it lies in the convex hull of the points $x+h_S$, all in $[1,P]$; hence $p_j(x;\mathbf h)\ge\frac{k!}{(k-j)!}\ge1$, and for $j<k$, $p_j$ is strictly increasing in $x$. For $j=k-1$, $p_{k-1}(x;\mathbf h)=k!\,(x+\frac{h_1+\cdots+h_{k-1}}2)$.

\begin{theorem}[Weyl's inequality]\label{thm:weylk}
Let $k\ge2$ and $\varepsilon>0$. If $|\alpha-\frac aq|\le\frac1{q^2}$ with $(a,q)=1$, then
\[
|f(\alpha)|\ll_{k,\varepsilon}P^{1+\varepsilon}\Bigl(\frac1q+\frac1P+\frac q{P^k}\Bigr)^{2^{1-k}}.
\]
\end{theorem}
\begin{proof}
We may assume $q\le P^k$, since otherwise the bound exceeds $P$. Apply Lemma~\ref{lem:weyldiff} with $j=k-1$ and $I=[1,P]$. By ($\dagger$), $\Delta_{\mathbf h}(\alpha x^k)$ is linear in $x$ with slope $\alpha k!h_1\cdots h_{k-1}$, so the reminder bounds each inner sum by $\min(P+1,\frac1{2\|\alpha k!h_1\cdots h_{k-1}\|})$, and
\[
|f(\alpha)|^{2^{k-1}}\le(2P)^{2^{k-1}-k}\sum_{\mathbf h\in H^{k-1}}\min\Bigl(2P,\frac1{\|\alpha k!h_1\cdots h_{k-1}\|}\Bigr).
\]
The $\mathbf h$ with some $h_i=0$ contribute at most $(k-1)(2P)^{k-2}\cdot2P\ll P^{k-1}$. For the others, $m=k!|h_1\cdots h_{k-1}|$ lies in $[1,k!P^{k-1}]$, and a given $m$ comes from at most $2^{k-1}\tau(m)^{k-1}\ll P^\varepsilon$ vectors $\mathbf h$, since each $|h_i|$ divides $m$. With $X=k!P^{k-1}$, the second reminder bounds the sum over $m$ by
\[
\Bigl(\frac Xq+1\Bigr)(8P+2q+2q\log q)\ll\Bigl(\frac{P^k}q+P^{k-1}+P+q\Bigr)\log P\ll P^k\Bigl(\frac1q+\frac1P+\frac q{P^k}\Bigr)\log P.
\]
Altogether $|f(\alpha)|^{2^{k-1}}\ll P^{2^{k-1}-k}\cdot P^{k+2\varepsilon}(\frac1q+\frac1P+\frac q{P^k})$; take the $2^{k-1}$-th root and rename $\varepsilon$.
\end{proof}

For $k=2$ the saving is a square root. The exponent $2^{1-k}$ is the price of $k-1$ applications of Cauchy--Schwarz; Vinogradov's mean value theorem, quoted in Chapter~\ref{ch:expsums}, improves it to about $\frac1{k^2}$, but we shall not need that.

\begin{theorem}[Hua, 1938]\label{thm:hua}
For $1\le j\le k$ and $\varepsilon>0$, $\displaystyle\int_0^1|f(\alpha)|^{2^j}\dd\alpha\ll_{k,\varepsilon}P^{2^j-j+\varepsilon}$.
\end{theorem}
\begin{proof}
For $j=1$, Parseval gives $\int_0^1|f|^2=\lfloor P\rfloor\le P$. Suppose the bound holds for some $j\le k-1$. Expanding,
\[
|f(\alpha)|^{2^j}=f(\alpha)^{2^{j-1}}\overline{f(\alpha)}^{2^{j-1}}=\sum_m\rho(m)e(\alpha m),
\]
where $\rho(m)\ge0$ counts the solutions of $y_1^k+\cdots+y_{2^{j-1}}^k-z_1^k-\cdots-z_{2^{j-1}}^k=m$ in integers of $[1,P]$; and $\rho(m)\le\rho(0)=\int_0^1|f|^{2^j}$ by Cauchy--Schwarz, $\sum_m\rho(m)=\lfloor P\rfloor^{2^j}$. Lemma~\ref{lem:weyldiff} and ($\dagger$) give
\[
|f(\alpha)|^{2^j}\le(2P)^{2^j-j-1}\sum_{\mathbf h\in H^j}\sum_{x\in I_{\mathbf h}}e\bigl(\alpha D(x;\mathbf h)\bigr),\qquad D(x;\mathbf h)=h_1\cdots h_jp_j(x;\mathbf h),
\]
with a non-negative right side. Multiplying by $|f(\alpha)|^{2^j}$ and integrating, since $\int_0^1e(\alpha D)|f|^{2^j}=\rho(-D)$,
\[
\int_0^1|f|^{2^{j+1}}\le(2P)^{2^j-j-1}\sum_{\mathbf h\in H^j}\sum_{x\in I_{\mathbf h}}\rho(-D(x;\mathbf h)).
\]
If some $h_i=0$, then $D=0$; there are at most $j(2P)^{j-1}P$ such pairs $(\mathbf h,x)$, each contributing $\rho(0)\ll P^{2^j-j+\varepsilon}$, in total $O(P^{2^j+\varepsilon})$. If no $h_i$ is $0$, then $D\ne0$ and $|D|\le k!P^k$. A given value $m=D$ determines the $h_i$ up to $(2\tau(|m|))^j\ll P^\varepsilon$ choices, since each divides $m$; then $p_j=\frac m{h_1\cdots h_j}$, which determines $x$ because $p_j$ is strictly increasing in $x$. So these pairs contribute at most $O(P^\varepsilon)\sum_m\rho(-m)\ll P^{2^j+\varepsilon}$. Hence $\int_0^1|f|^{2^{j+1}}\ll P^{2^j-j-1}\cdot P^{2^j+\varepsilon}=P^{2^{j+1}-(j+1)+\varepsilon}$.
\end{proof}

For $j=k$ the lemma says that $\int|f|^{2^k}\ll P^{2^k-k+\varepsilon}$, which is the number of trivial solutions $P^{2^{k-1}}$ only when $k=2$, but is always within $P^\varepsilon$ of the lower bound $P^{2^k-k}$: the $P^{2^k}$ tuples of $2^k$ variables produce values in a range of length $O(P^k)$, so some value is hit $P^{2^k-k}$ times. Hua's lemma says that, on average, the sums of $2^{k-1}$ $k$-th powers are as well spread as they can be.

\section{Waring's problem}

Let $r_{s,k}(n)$ be the number of $(x_1,\dots,x_s)$ in positive integers with $x_1^k+\cdots+x_s^k=n$. Take $P=n^{1/k}$; then every solution has $x_i\le P$ and $r_{s,k}(n)=\int_0^1f(\alpha)^se(-n\alpha)\dd\alpha$. The major arc approximation involves the complete sums
\[
S(q,a)=\sum_{r=1}^qe\Bigl(\frac{ar^k}q\Bigr),\qquad v(\beta)=\int_0^Pe(\beta t^k)\dd t.
\]

\begin{lemma}\label{lem:waringmajor}
If $1\le q\le P$ and $\alpha=\frac aq+\beta$, then $f(\alpha)=q^{-1}S(q,a)v(\beta)+O_k(q(1+P^k|\beta|))$.
\end{lemma}
\begin{proof}
Since $x^k\equiv r^k\pmod q$ when $x\equiv r$,
\[
f(\alpha)=\sum_{r=1}^qe\Bigl(\frac{ar^k}q\Bigr)\sum_{\substack{x\le P\\x\equiv r}}g(x),\qquad g(t)=e(\beta t^k).
\]
 For fixed $r$, write $x=r+qm$, $0\le m\le M$. Each term $g(r+qm)$ differs from $\frac1q\int_{r+qm}^{r+q(m+1)}g$ by at most $\int_{r+qm}^{r+q(m+1)}|g'|$, and $\frac1q\int_{r}^{r+q(M+1)}g$ differs from $\frac1qv(\beta)$ by at most $2$, the two pieces $[0,r]$ and $[P,r+q(M+1)]$ having length at most $q$. With $|g'(t)|=2\pi k|\beta|t^{k-1}$ and $q\le P$, $\int_0^{2P}|g'|=2\pi|\beta|(2P)^k$. So the inner sum is $\frac{v(\beta)}q+O(1+P^k|\beta|)$; multiply by the $q$ roots of unity and add.
\end{proof}

\begin{lemma}\label{lem:singint}
(a) $|v(\beta)|\le\min(P,2|\beta|^{-1/k})$.
(b) If $s\ge2k$, then $\displaystyle\int_\R v(\beta)^se(-n\beta)\dd\beta=\frac{\Gamma(1+\frac1k)^s}{\Gamma(\frac sk)}\,n^{\frac sk-1}$.
\end{lemma}
\begin{proof}
Substituting $u=t^k$, $v(\beta)=\int_0^ng(u)e(\beta u)\dd u$ with $g(u)=\frac1ku^{1/k-1}$ on $\left]0,n\right]$ and $g=0$ elsewhere.
(a) $|v|\le P$ is clear. For $|\beta|\ge\frac1n$ let $U=\frac1{|\beta|}$; then $\int_0^Ug=U^{1/k}$. On $[U,n]$, $g$ is positive and decreasing, and $E(u)=\int_U^ue(\beta w)\dd w$ has $|E|\le\frac1{\pi|\beta|}$; integrating by parts, $|\int_U^nge(\beta u)\dd u|\le\frac{2g(U)}{\pi|\beta|}=\frac{2U^{1/k}}{k\pi}$. So $|v(\beta)|\le(1+\frac2{k\pi})U^{1/k}\le2|\beta|^{-1/k}$; for $|\beta|<\frac1n$, $P<|\beta|^{-1/k}$.
(b) By (a), $|v|^s$ is integrable when $s>k$. Let $g^{*j}$ be the $j$-fold convolution of $g$. The Beta integral gives, by induction on $j$, $\int_0^m\frac{u^{a-1}}{\Gamma(a)}\frac{(m-u)^{b-1}}{\Gamma(b)}\dd u=\frac{m^{a+b-1}}{\Gamma(a+b)}$, hence
\[
g^{*j}(m)=\frac{\Gamma(\frac1k)^j}{k^j\Gamma(\frac jk)}m^{\frac jk-1}=\frac{\Gamma(1+\frac1k)^j}{\Gamma(\frac jk)}m^{\frac jk-1}\qquad(0<m\le n),
\]
because for $m\le n$ the cut-off at $n$ plays no role; for $m>n$, $g^{*j}(m)$ is at most the same expression. In particular $g^{*k}$ is bounded, so $g^{*2k}=g^{*k}*g^{*k}$ is continuous (a convolution of two functions in $L^2$), and so is $g^{*s}=g^{*2k}*g^{*(s-2k)}$ when $s>2k$, a convolution of a bounded uniformly continuous function with an integrable one. Its Fourier transform is $\hat g(\xi)^s=v(-\xi)^s$, which is integrable. By Fourier inversion at the point $n$, which is allowed because $g^{*s}$ is continuous, $g^{*s}(n)=\int v(-\xi)^se(n\xi)\dd\xi=\int v(\beta)^se(-n\beta)\dd\beta$.
\end{proof}

The condition $s\ge2k$ in (b) only serves to make the continuity easy; it holds in the theorem below, since $2^k+1\ge2k$.

\begin{theorem}[Hardy and Littlewood, 1920; Hua, 1938]\label{thm:waring}
Let $k\ge2$ and $s\ge2^k+1$. The series
\[
\mathfrak S_{s,k}(n)=\sum_{q\ge1}\sum_{\substack{a\bmod q\\(a,q)=1}}\Bigl(\frac{S(q,a)}q\Bigr)^se\Bigl(-\frac{an}q\Bigr)
\]
converges absolutely, uniformly in $n$, and there is $\delta=\delta(k)>0$ such that
\[
r_{s,k}(n)=\frac{\Gamma(1+\frac1k)^s}{\Gamma(\frac sk)}\,\mathfrak S_{s,k}(n)\,n^{\frac sk-1}+O_{s,k}\bigl(n^{\frac sk-1-\delta}\bigr).
\]
\end{theorem}
\begin{proof}
Fix $\varepsilon>0$ small in terms of $s$ and $k$, and let $\nu=\frac1{10}$.

\emph{Complete sums.} For $(a,q)=1$, $S(q,a)$ is the Weyl sum of length $q$ at $\alpha=\frac aq$, so Theorem~\ref{thm:weylk} with $P=q$ gives $|S(q,a)|\ll q^{1+\varepsilon}(\frac2q+q^{1-k})^{2^{1-k}}\ll q^{1-2^{1-k}+\varepsilon}$. As $s2^{1-k}\ge2+2^{1-k}$, the inner sum in $\mathfrak S_{s,k}$ is $O(q\cdot q^{-s2^{1-k}+s\varepsilon})=O(q^{-1-2^{1-k}+s\varepsilon})$, so the series converges absolutely, and its tail beyond $Q$ is $O(Q^{-2^{1-k}+s\varepsilon})$.

\emph{Arcs.} For $q\le P^\nu$ and $(a,q)=1$ let $\mathfrak M(q,a)=\{\alpha:|\alpha-\frac aq|\le P^{\nu-k}\}$; they are disjoint for $P$ large, as $P^{-2\nu}>2P^{\nu-k}$, and their union $\mathfrak M$ has measure at most $2P^{3\nu-k}$. On the complement $\mathfrak m$, Lemma~\ref{lem:dirichlet} with $Q=P^{k-\nu}$ gives $q\le Q$ with $|\alpha-\frac aq|\le\frac1{qQ}\le\frac1{q^2}$, and $q>P^\nu$ (otherwise $\alpha\in\mathfrak M$). Theorem~\ref{thm:weylk} then gives $|f(\alpha)|\ll P^{1+\varepsilon}(3P^{-\nu})^{2^{1-k}}$.

\emph{Minor arcs.} By Hua's lemma with $j=k$,
\[
\int_{\mathfrak m}|f|^s\le\sup_{\mathfrak m}|f|^{s-2^k}\int_0^1|f|^{2^k}\ll P^{(s-2^k)(1-\nu2^{1-k}+\varepsilon)}P^{2^k-k+\varepsilon}\le P^{s-k-\nu2^{1-k}+s\varepsilon}.
\]

\emph{Major arcs.} On $\mathfrak M(q,a)$, with $V=q^{-1}S(q,a)v(\beta)$, Lemma~\ref{lem:waringmajor} gives $|f-V|\ll P^\nu(1+P^\nu)\ll P^{2\nu}$, and $|f|,|V|\le P$, so $|f^s-V^s|\le sP^{s-1}|f-V|\ll P^{s-1+2\nu}$. Over $\mathfrak M$ this costs $O(P^{s-k-1+5\nu})$. The main part is
\[
\sum_{q\le P^\nu}\sum_{(a,q)=1}\Bigl(\frac{S(q,a)}q\Bigr)^se\Bigl(-\frac{an}q\Bigr)\int_{|\beta|\le P^{\nu-k}}v(\beta)^se(-n\beta)\dd\beta.
\]
By Lemma~\ref{lem:singint}(a), the integral over $|\beta|>P^{\nu-k}$ is at most $2\int_{P^{\nu-k}}^\infty(2\beta^{-1/k})^s\dd\beta\ll P^{(\nu-k)(1-s/k)}=P^{s-k-\nu(s/k-1)}$, so the integral is $J(n)+O(P^{s-k-\nu(s-k)/k})$ with $J(n)=\frac{\Gamma(1+1/k)^s}{\Gamma(s/k)}P^{s-k}$ by Lemma~\ref{lem:singint}(b). The sum over $q\le P^\nu$ is bounded, and differs from $\mathfrak S_{s,k}(n)$ by $O(P^{-\nu(2^{1-k}-s\varepsilon)})$.

\emph{Conclusion.} Every error is $O(P^{s-k-\delta'})$ for some $\delta'>0$ once $\varepsilon$ is small: the exponents saved are $\nu2^{1-k}-s\varepsilon$, $1-5\nu$, $\nu\frac{s-k}k$ and $\nu(2^{1-k}-s\varepsilon)$. Since $P^{s-k}=n^{s/k-1}$, the theorem holds with $\delta=\frac{\delta'}k$.
\end{proof}

Figure~\ref{fig:fivesquares} tests the formula on five squares, the case $k=2$, $s=2^k+1$. The singular series, computed from its first $200$ terms, predicts the size of $r_{5,2}(n)$ from one $n$ to the next, and the scatter around a straight line is small already at $n=4000$.

\section{The singular series}

The main term of Theorem~\ref{thm:waring} is useful only if $\mathfrak S_{s,k}(n)$ is not small. It can vanish. Fourth powers are $0$ or $1$ modulo $16$, so a sum of at most $14$ fourth powers is never $15$ modulo $16$; for $s\le14$ and $n\equiv15\pmod{16}$ there is no representation at all, not even modulo $16$. The singular series knows this. It is a product over the primes, and the factor at $p$ measures how often $n$ is a sum of $s$ $k$-th powers modulo high powers of $p$. This section proves the product formula, and then shows that with $s\ge4k$ every factor is positive and the product is bounded below.

Write $A(q,n)=\sum_{a\bmod q,\,(a,q)=1}\bigl(\frac{S(q,a)}q\bigr)^se\bigl(-\frac{an}q\bigr)$, so that $\mathfrak S_{s,k}(n)=\sum_{q\ge1}A(q,n)$ and $A(1,n)=1$. The proof of Theorem~\ref{thm:waring} showed, for $s\ge2^k+1$ and every $\varepsilon>0$, that
\[
|A(q,n)|\le C(s,k,\varepsilon)\,q^{-1-2^{1-k}+s\varepsilon}\qquad\text{uniformly in }n.\tag{$\ddagger$}
\]
For a prime power $p^h$ let $M_n(p^h)$ be the number of $(x_1,\dots,x_s)$ modulo $p^h$ with $x_1^k+\cdots+x_s^k\equiv n\pmod{p^h}$.

\begin{lemma}\label{lem:amult}
If $q=q_1q_2$ with $(q_1,q_2)=1$, then $A(q,n)=A(q_1,n)A(q_2,n)$.
\end{lemma}
\begin{proof}
By the Chinese remainder theorem, as $a_1$ runs over the units modulo $q_1$ and $a_2$ over the units modulo $q_2$, $a=a_1q_2+a_2q_1$ runs once over the units modulo $q$, and $\frac aq=\frac{a_1}{q_1}+\frac{a_2}{q_2}$. In the same way $x=x_1q_2+x_2q_1$ runs once over the residues modulo $q$. Expanding, $(x_1q_2+x_2q_1)^k=(x_1q_2)^k+(x_2q_1)^k+q_1q_2m$ with $m$ an integer, so modulo $1$
\[
\frac{ax^k}q\equiv\frac{a_1(x_1q_2)^k}{q_1}+\frac{a_2(x_2q_1)^k}{q_2},
\]
because each cross term is an integer. As $x_1$ runs modulo $q_1$, so does $x_1q_2$. Hence $S(q,a)=S(q_1,a_1)S(q_2,a_2)$. Also $e(-\frac{an}q)=e(-\frac{a_1n}{q_1})e(-\frac{a_2n}{q_2})$, and the double sum over $a_1,a_2$ factorizes.
\end{proof}

\begin{lemma}\label{lem:localcount}
For every prime $p$ and $h\ge1$, $\displaystyle\sum_{j=0}^hA(p^j,n)=p^{h(1-s)}M_n(p^h)$.
\end{lemma}
\begin{proof}
Orthogonality of the characters $b\mapsto e(\frac{bm}{p^h})$ gives
\[
M_n(p^h)=\frac1{p^h}\sum_{b\bmod p^h}\ \sum_{x_1,\dots,x_s\bmod p^h}e\Bigl(\frac{b(x_1^k+\cdots+x_s^k-n)}{p^h}\Bigr)=\frac1{p^h}\sum_{b\bmod p^h}S(p^h,b)^se\Bigl(-\frac{bn}{p^h}\Bigr).
\]
Group the $b$ by their gcd with $p^h$: $b=p^{h-j}a$ with $0\le j\le h$ and $a$ a unit modulo $p^j$ (for $j=0$, $b=0$ and $a$ runs over the single class modulo $1$). Each residue modulo $p^j$ is hit $p^{h-j}$ times by $x$ modulo $p^h$, so $S(p^h,p^{h-j}a)=p^{h-j}S(p^j,a)$. Therefore
\[
M_n(p^h)=\frac1{p^h}\sum_{j=0}^hp^{(h-j)s}\sum_{a}S(p^j,a)^se\Bigl(-\frac{an}{p^j}\Bigr)=p^{h(s-1)}\sum_{j=0}^hA(p^j,n).\qedhere
\]
\end{proof}

So the partial sums of $\chi_p(n)=\sum_{j\ge0}A(p^j,n)$ are densities: $p^{h(1-s)}M_n(p^h)$ is the number of solutions modulo $p^h$ divided by the number $p^{h(s-1)}$ that a random value of $x_1^k+\cdots+x_s^k$ would give. To show that $\chi_p(n)>0$ we produce many solutions modulo $p^h$, by finding one modulo a fixed power of $p$ and lifting it. Write $k=p^\tau k_0$ with $p\nmid k_0$, and let $\gamma=\tau+1$ if $p$ is odd, $\gamma=\tau+2$ if $p=2$.

\begin{lemma}[lifting $k$-th powers]\label{lem:henselk}
Let $h\ge\gamma$, let $x$ be an integer prime to $p$, and let $b\equiv x^k\pmod{p^h}$. Then there is an integer $x'\equiv x\pmod{p^{h-\tau}}$ with $x'^k\equiv b\pmod{p^{h+1}}$. Consequently, if $b\equiv y^k\pmod{p^\gamma}$ with $p\nmid y$, then $b$ is a $k$-th power modulo every $p^h$.
\end{lemma}
\begin{proof}
Write $b=x^k+cp^h$ and try $x'=x+tp^{h-\tau}$ with $t$ an integer. By the binomial theorem,
\[
x'^k=x^k+kx^{k-1}tp^{h-\tau}+\sum_{j=2}^k\binom kjx^{k-j}t^jp^{j(h-\tau)}.
\]
The linear term is $k_0x^{k-1}t\,p^h$. For $j\ge2$, $\binom kj=\frac kj\binom{k-1}{j-1}$ shows that $p^{\tau-v}$ divides $\binom kj$, where $p^v$ is the exact power of $p$ dividing $j$. So the $j$-th term is divisible by $p^{h+1}$ as soon as $\tau-v+j(h-\tau)\ge h+1$, that is $(j-1)(h-\tau)\ge1+v$. For $p$ odd, $h-\tau\ge1$, and $1+v\le j-1$: this is $1\le1$ for $j=2$, and for $j\ge3$ it holds because $3^v\le j$ gives $v\le j-2$. For $p=2$, $h-\tau\ge2$, and $1+v\le2(j-1)$ because $2^v\le j$. So $x'^k\equiv x^k+k_0x^{k-1}t\,p^h\pmod{p^{h+1}}$, and we need $k_0x^{k-1}t\equiv c\pmod p$. Since $p\nmid k_0x$, there is such a $t$. The last sentence follows by induction on $h$, starting from $x=y$ at $h=\gamma$.
\end{proof}

\begin{lemma}\label{lem:localsol}
Every residue class modulo $p^\gamma$ is a sum of at most $4k-1$ numbers of the form $u^k$ with $p\nmid u$.
\end{lemma}
\begin{proof}
Let $G=\Z/p^\gamma\Z$, $U$ its group of units, $H=\{u^k:u\in U\}$, and let $T_j$ be the set of sums of at most $j$ elements of $H$, with $T_0=\{0\}$. Since $1=1^k\in H$, $T_j$ contains $0,1,\dots,j$.

\emph{The case $p=2$.} Here $p^\gamma=2^{\tau+2}=4\cdot2^\tau\le4k$, so the classes $0,1,\dots,4k-1$ already cover $G$.

\emph{The case $p$ odd.} Here $\gamma=\tau+1$. First, the only element of $H$ that is $\equiv1\pmod p$ is $1$. Indeed, $U$ has $p^\tau(p-1)$ elements and its subgroup $K=\{u\equiv1\pmod p\}$ has $p^\tau$. If $h=u^k\in H\cap K$, then $h^{p-1}=(u^{p^\tau(p-1)})^{k_0}=1$ by Lagrange's theorem in $U$, and $h^{p^\tau}=1$ by Lagrange's theorem in $K$; as $\gcd(p-1,p^\tau)=1$, $h=1$. Two consequences follow.

(i) Reduction modulo $p$ is injective on $H$: if $h_1\equiv h_2\pmod p$, then $h_1h_2^{-1}\in H\cap K$. Its image is the set of $k$-th powers in $\mathbb F_p^\times$, and $x^k=(x^{p^\tau})^{k_0}=x^{k_0}$ there by Fermat's little theorem. Each value of $x\mapsto x^{k_0}$ is taken at most $k_0$ times, since a polynomial of degree $k_0$ has at most $k_0$ roots. So $|H|\ge\frac{p-1}{k_0}$.

(ii) For $g\ne0$ in $G$, the elements $hg$, $h\in H$, are distinct. If $hg=g$, write $g=p^eg'$ with $e<\gamma$ and $p\nmid g'$; then $p^{\gamma-e}$ divides $h-1$, so $h\equiv1\pmod p$ and $h=1$.

Now $T_{j+1}=T_j\cup(T_j+H)$, and $hT_j=T_j$ for $h\in H$, since $H$ is closed under multiplication. If $T_{j+1}=T_j$, then $T_j+1\subseteq T_j$, so $T_j$ contains $0,1,2,\dots$ and $T_j=G$. Otherwise $T_{j+1}\setminus T_j$ is non-empty, does not contain $0$, and is stable under multiplication by $H$, so by (ii) it has at least $|H|$ elements. By induction, $|T_j|\ge\min\bigl(p^\gamma,1+j\frac{p-1}{k_0}\bigr)$. The second number reaches $p^{\tau+1}$ once $j\ge k_0(1+p+\cdots+p^\tau)$, and
\[
k_0(1+p+\cdots+p^\tau)\le k_0p^\tau\Bigl(1+\frac1p+\frac1{p^2}+\cdots\Bigr)=k\,\frac p{p-1}\le\frac32k\le4k-1.\qedhere
\]
\end{proof}

The bound $4k$ cannot be lowered in general. For $k=4$ and $p=2$, $\gamma=4$, every odd fourth power is $1$ modulo $16$ and every even one is $0$, so the class $15$ needs fifteen odd fourth powers (Exercise~\ref{ex:sharp4k}).

\begin{theorem}[Hardy and Littlewood]\label{thm:singpos}
Let $k\ge2$ and $s\ge\max(4k,2^k+1)$. There is $c(s,k)>0$ such that $\mathfrak S_{s,k}(n)\ge c(s,k)$ for every $n\ge1$.
\end{theorem}
\begin{proof}
\emph{Each factor is positive.} Fix $p$ and $\gamma$ as above. By Lemma~\ref{lem:localsol} applied to the class of $n-1$, there are units $u_1,\dots,u_j$ with $j\le4k-1\le s-1$ and $n\equiv1+u_1^k+\cdots+u_j^k\pmod{p^\gamma}$. So $(1,u_1,\dots,u_j,0,\dots,0)$ is a solution modulo $p^\gamma$ whose first coordinate is a unit. Let $h\ge\gamma$. Choose $x_2,\dots,x_s$ modulo $p^h$ congruent modulo $p^\gamma$ to $u_1,\dots,u_j,0,\dots,0$; there are $p^{(h-\gamma)(s-1)}$ choices. For each, $b=n-x_2^k-\cdots-x_s^k\equiv1=1^k\pmod{p^\gamma}$, so by Lemma~\ref{lem:henselk} there is $x_1$ with $x_1^k\equiv b\pmod{p^h}$. Hence $M_n(p^h)\ge p^{(h-\gamma)(s-1)}$, and by Lemma~\ref{lem:localcount}
\[
\sum_{j=0}^hA(p^j,n)=p^{h(1-s)}M_n(p^h)\ge p^{-\gamma(s-1)}.
\]
The series $\chi_p(n)=\sum_{j\ge0}A(p^j,n)$ converges by ($\ddagger$), so $\chi_p(n)\ge p^{-\gamma(s-1)}$.

\emph{The product formula.} For $y\ge2$, multiply out the finite product $\prod_{p\le y}\chi_p(n)$ of absolutely convergent series. By Lemma~\ref{lem:amult} and unique factorization, it is $\sum A(q,n)$ over the $q$ whose prime factors are all at most $y$. This differs from $\mathfrak S_{s,k}(n)$ by at most $\sum_{q>y}|A(q,n)|$, which tends to $0$. So $\mathfrak S_{s,k}(n)=\lim_{y\to\infty}\prod_{p\le y}\chi_p(n)$.

\emph{The large primes.} Fix $\varepsilon=\frac{2^{-k}}s$, so that $\eta=2^{1-k}-s\varepsilon>0$ in ($\ddagger$). Let $E_p=\sum_{j\ge1}|A(p^j,n)|$. Then $\sum_{p>p_0}E_p\le\sum_{q>p_0}Cq^{-1-\eta}\le\frac12$ for some $p_0=p_0(s,k)$ that does not depend on $n$. For $p>p_0$, $\chi_p(n)\ge1-E_p\ge\frac12$, and for every $y$,
\[
\prod_{p_0<p\le y}\chi_p(n)\ge\prod_{p_0<p\le y}(1-E_p)\ge1-\sum_{p>p_0}E_p\ge\frac12,
\]
using $(1-a)(1-b)\ge1-a-b$ for $a,b\in[0,1]$, repeatedly.

\emph{Conclusion.} For $p\le p_0$, $\gamma\le\tau+2\le\log_2k+2$. So $\mathfrak S_{s,k}(n)\ge\frac12\prod_{p\le p_0}p^{-(\log_2k+2)(s-1)}$, a positive number that depends only on $s$ and $k$.
\end{proof}

\begin{corollary}[Hardy and Littlewood; Hua, 1938]\label{cor:gk}
For $k\ge2$, every sufficiently large integer is a sum of $\max(4k,2^k+1)$ positive $k$-th powers. In particular $G(k)\le2^k+1$ for $k\ge4$, $G(3)\le12$ and $G(2)\le8$.
\end{corollary}
\begin{proof}
Let $s=\max(4k,2^k+1)$. Theorems~\ref{thm:waring} and~\ref{thm:singpos} give $r_{s,k}(n)\ge c'\,n^{s/k-1}-O(n^{s/k-1-\delta})$ with $c'>0$, which is positive for $n$ large. For $k\ge4$, $2^k+1\ge4k$.
\end{proof}

The small cases are far from the truth: Lagrange gives $G(2)=4$, and Linnik proved $G(3)\le7$ in 1943. The point of the corollary is that one argument works for every $k$. Refinements of Vinogradov's methods bring $G(k)$ down to about $k\log k$; Wooley proved $G(k)\le k(\log k+\log\log k+O(1))$ in 1992. The local obstruction is real: a number $31\cdot16^m$ needs sixteen fourth powers, so $G(4)\ge16$, and Davenport proved $G(4)=16$ in 1939.

\begin{figure}[ht]
\centering
\fig{waring-fivesquares}
\caption{Sums of five squares ($k=2$, $s=5=2^k+1$), for $n=2000,2010,\dots,4000$. Horizontally, the singular series $\mathfrak S_{5,2}(n)$, summed over $q\le200$; vertically, $r_{5,2}(n)$, computed exactly, divided by $\frac{\Gamma(3/2)^5}{\Gamma(5/2)}n^{3/2}$. The correlation is $0.993$. The points lie along a line slightly below the diagonal (the average ratio is $0.93$), because of a lower-order term of relative size about $n^{-1/2}$, coming from the representations that would use $0$.}
\label{fig:fivesquares}
\end{figure}

\begin{beauty}
A count of solutions is written as an integral over the circle. Near the fractions with small denominators the integrand factorizes into a \emph{local} part, a complete exponential sum modulo $q$, and an \emph{archimedean} part, an integral over the reals. Summing over all $q$ assembles the local parts into a product over all primes. The main term is then a product of densities, one for each completion of $\Q$: the singular integral for $\R$, the factor $\chi_p$ for each $\Q_p$. Everything else, the minor arcs, is noise, and the whole difficulty is to prove that it is noise. Chapter~\ref{ch:threeprimes} runs the same argument for three primes.
\end{beauty}

\section{Exercises}

\begin{exercise}[$\star$]\label{ex:gausssq}
For $k=2$ and $(a,q)=1$, show that $|S(q,a)|^2\le2q$. Deduce directly, without Weyl's inequality, that $\mathfrak S_{s,2}(n)$ converges absolutely for $s\ge5$.
\end{exercise}
\begin{exercise}[$\star\star$]\label{ex:density}
Use Hua's lemma and Cauchy--Schwarz to show that at least $X^{1-\varepsilon}$ of the integers $n\le X$ are sums of $2^{k-1}$ positive $k$-th powers, for $X\ge X_0(k,\varepsilon)$. What does it say for four cubes?
\end{exercise}
\begin{exercise}[$\star$]\label{ex:g4}
(a) Show that $79$ is not a sum of fewer than nineteen positive fourth powers. (b) Find all ways of writing $79$ as a sum of nineteen fourth powers, up to order.
\end{exercise}
\begin{exercise}[$\star\star$]\label{ex:sharp4k}
(a) Show that $x^4\equiv0$ or $1\pmod{16}$ for every integer $x$, and deduce that a number $n\equiv15\pmod{16}$ is not a sum of fewer than fifteen fourth powers. Why does this show that $4k-1$ in Lemma~\ref{lem:localsol} cannot be lowered for $k=4$, $p=2$?
(b) Show that $31\cdot16^m$ is not a sum of fifteen or fewer fourth powers, for every $m\ge0$. Deduce $G(4)\ge16$.
\end{exercise}
\begin{exercise}[$\star\star\star$]\label{ex:foursq}
Let $k=2$ and let $p$ be an odd prime.
(a) Show that $S(p,a)=\bigl(\frac ap\bigr)g_p$ for $p\nmid a$, where $g_p$ is the Gauss sum of the reminder in Chapter~\ref{ch:expsums}, and deduce $|A(p,n)|\le p^{1-s/2}$.
(b) Take $s=4$ and $p\nmid n$. Compute $A(p,n)$ exactly, and deduce from Lemma~\ref{lem:localcount} that the congruence $x_1^2+x_2^2+x_3^2+x_4^2\equiv n\pmod p$ has exactly $p^3-p$ solutions. Check this by hand for $p=3$, $n=1$.
\end{exercise}
\section{Hints}

\begin{hint}{ex:gausssq}
Expand $|S(q,a)|^2$, put $x=y+h$, and sum over $y$ first: $x^2-y^2=h^2+2hy$.
\end{hint}
\begin{hint}{ex:density}
$\sum_nr_{s,k}(n)=\lfloor P\rfloor^s$ and $\sum_nr_{s,k}(n)^2=\int_0^1|f|^{2s}$; take $2s=2^k$.
\end{hint}
\begin{hint}{ex:g4}
Only $1$ and $16$ are available. Count how many $16$'s can be used.
\end{hint}
\begin{hint}{ex:sharp4k}
For (b), if $m\ge1$ and fifteen fourth powers sum to a multiple of $16$, how many of them can be odd?
\end{hint}
\begin{hint}{ex:foursq}
For (a), $\sum_xe(\frac{ax^2}p)=\sum_y\bigl(1+(\frac yp)\bigr)e(\frac{ay}p)$. For (b), $g_p^4=p^2$.
\end{hint}
\section{Solutions}

\begin{solution}{ex:gausssq}
\idea Squaring a Gauss sum linearizes the phase in one of the two variables.
\writeup $|S(q,a)|^2=\sum_{x,y\bmod q}e(\frac{a(x^2-y^2)}q)$. With $x=y+h$, $x^2-y^2=h^2+2hy$, so
\[
|S(q,a)|^2=\sum_{h\bmod q}e\Bigl(\frac{ah^2}q\Bigr)\sum_{y\bmod q}e\Bigl(\frac{2ahy}q\Bigr).
\]
The inner sum is $q$ if $q\mid2ah$, that is $q\mid2h$ since $(a,q)=1$, and $0$ otherwise. At most two classes $h$ modulo $q$ have $q\mid2h$, so $|S(q,a)|^2\le2q$. Then the $q$-th term of $\mathfrak S_{s,2}$ is at most $q(\frac{2}{q})^{s/2}=2^{s/2}q^{1-s/2}$, and $\sum_qq^{1-s/2}$ converges for $s\ge5$.
\end{solution}
\begin{solution}{ex:density}
\idea If the values were concentrated on few integers, the number of coincidences, which is the $2s$-th moment, would be large.
\writeup Let $s=2^{k-1}$, $P=X^{1/k}$, and let $Z$ be the number of $n\le X$ with $r_{s,k}(n)>0$; all values $\sum x_i^k$ with $x_i\le P/s^{1/k}$ are at most $X$, so with $P'=P/s^{1/k}$ and $r'$ counting those, Cauchy--Schwarz gives
\[
\lfloor P'\rfloor^{2s}=\Bigl(\sum_nr'(n)\Bigr)^2\le Z\sum_nr'(n)^2=Z\int_0^1|f_{P'}|^{2^k}\ll ZP'^{2^k-k+\varepsilon}
\]
by Hua's lemma with $j=k$ ($2s=2^k$). So $Z\gg P'^{k-\varepsilon}\gg X^{1-\varepsilon/k}$. For cubes, $s=4$: almost every integer, in the sense that $X^{1-\varepsilon}$ of those up to $X$, is a sum of four positive cubes. Much more is true: Davenport proved in 1939, with a refined circle method, that almost all integers are sums of four positive cubes.
\end{solution}
\begin{solution}{ex:g4}
\idea With so few powers available, a representation is just a choice of how many $16$'s to use.
\writeup (a) Since $3^4=81>79$, the only positive fourth powers at most $79$ are $1$ and $16$. A representation with $a$ sixteens and $b$ ones has $16a+b=79$, so $a\le4$ and the number of terms is $a+b=79-15a\ge79-60=19$. (b) Equality needs $a=4$, $b=15$: the only representation with nineteen terms is $4\cdot2^4+15\cdot1^4$.
\end{solution}
\begin{solution}{ex:sharp4k}
\idea Modulo $16$ a fourth power only says whether its base is odd.
\writeup (a) If $x$ is even, $16\mid x^4$. If $x$ is odd, $x^2=8t+1$ for some $t$ (since $x^2-1=(x-1)(x+1)$ is a product of consecutive even numbers), so $x^4=64t^2+16t+1\equiv1\pmod{16}$. A sum of $s$ fourth powers is therefore congruent to the number of odd bases, which is between $0$ and $s$. For $n\equiv15\pmod{16}$ we need at least fifteen odd bases. Modulo $p^\gamma=2^{2+2}=16$, the unit fourth powers are all $1$, so the class $15$ is a sum of unit fourth powers only with fifteen terms, and $4k-1=15$.
(b) For $m=0$: $31<81$, so as in Exercise~\ref{ex:g4} a representation is $16a+b=31$ with $a\le1$, and has $a+b=31-15a\ge16$ terms. For $m\ge1$, suppose $31\cdot16^m=\sum_{i\le15}x_i^4$. The number of odd $x_i$ is $\equiv0\pmod{16}$ and at most $15$, so it is $0$: every $x_i$ is even, and dividing by $16$ gives $31\cdot16^{m-1}=\sum(\frac{x_i}2)^4$ with fifteen terms. Repeating, we reach $31$ as a sum of at most fifteen fourth powers, which is impossible. So infinitely many integers need sixteen fourth powers, and $G(4)\ge16$.
\end{solution}
\begin{solution}{ex:foursq}
\idea For squares the complete sum is a Gauss sum, whose size is exactly $\sqrt p$.
\writeup (a) The number of $x$ modulo $p$ with $x^2\equiv y$ is $1+(\frac yp)$. So $S(p,a)=\sum_y(1+(\frac yp))e(\frac{ay}p)=\sum_y(\frac yp)e(\frac{ay}p)$, because $\sum_ye(\frac{ay}p)=0$ for $p\nmid a$. By the reminder of Chapter~\ref{ch:expsums}, this is $(\frac ap)g_p$, and $|g_p|^2=|g_p^2|=p$. Hence $|A(p,n)|\le(p-1)\bigl(\frac{\sqrt p}p\bigr)^s\le p^{1-s/2}$.
(b) For $s=4$, $(\frac{S(p,a)}p)^4=\frac{g_p^4}{p^4}=\frac{p^2}{p^4}$, since $g_p^2=\pm p$. So $A(p,n)=\frac1{p^2}\sum_{a=1}^{p-1}e(-\frac{an}p)=-\frac1{p^2}$ when $p\nmid n$. Lemma~\ref{lem:localcount} with $h=1$ gives $p^{-3}M_n(p)=1-\frac1{p^2}$, so $M_n(p)=p^3-p$. For $p=3$, $n=1$: squares modulo $3$ are $0$ (from $x=0$) and $1$ (from $x=\pm1$). We need the number of nonzero coordinates to be $\equiv1\pmod3$, that is one or four of them. One nonzero coordinate: $4\cdot2=8$ choices. Four: $2^4=16$. Total $24=27-3$.
\end{solution}
\section{What we proved, and what comes next}

The epigraph asked how many ways a large number has of being a sum of $s$ $k$-th powers. For $s\ge2^k+1$ the answer is a formula: $r_{s,k}(n)$ is the singular integral $\frac{\Gamma(1+1/k)^s}{\Gamma(s/k)}n^{s/k-1}$ times the singular series $\mathfrak S_{s,k}(n)$, up to a smaller power of $n$ (Theorem~\ref{thm:waring}). Two estimates carry the proof. Weyl's inequality (Theorem~\ref{thm:weylk}) makes the generating function small away from fractions with small denominators. Hua's lemma (Theorem~\ref{thm:hua}) bounds its mean value, so that the minor arcs cost less than the main term. The singular series is a product of local densities (Lemmas~\ref{lem:amult} and~\ref{lem:localcount}), and when $s\ge4k$ every density is positive because a solution modulo $p^\gamma$ with a unit coordinate lifts to every power of $p$ (Lemmas~\ref{lem:henselk} and~\ref{lem:localsol}). So $G(k)\le2^k+1$ for $k\ge4$ (Corollary~\ref{cor:gk}).

The next chapter asks the same question with primes in place of $k$-th powers: in how many ways is an odd number a sum of three primes? The circle is the same, and so is the shape of the answer, a singular series times a singular integral. Both estimates change. On the minor arcs Weyl's inequality is replaced by Vinogradov's bound for sums over primes from Chapter~\ref{ch:expsums}. On the major arcs the complete sums $S(q,a)$ are replaced by the primes in progressions modulo $q$, which is where Siegel's theorem, and its constant that nobody can compute, enter.

\section*{Sources and further reading}

\begin{itemize}
\item \emph{History.} E.~Waring, \emph{Meditationes Algebraicae}, Cambridge, 1770. D.~Hilbert, ``Beweis für die Darstellbarkeit der ganzen Zahlen durch eine feste Anzahl $n^{\text{ter}}$ Potenzen (Waringsches Problem)'', \emph{Math.\ Ann.}\ 67 (1909), 281--300. On Hua: J.~J.~O'Connor and E.~F.~Robertson, ``Hua Loo-Keng'', MacTutor History of Mathematics archive.
\item \emph{Later results.} Yu.~V.~Linnik, ``On the representation of large numbers as sums of seven cubes'', \emph{Mat.\ Sb.}\ 12 (1943), 218--224. H.~Davenport, ``On Waring's problem for fourth powers'', \emph{Ann.\ of Math.}\ 40 (1939), 731--747. T.~D.~Wooley, ``Large improvements in Waring's problem'', \emph{Ann.\ of Math.}\ 135 (1992), 131--164. G.~H.~Hardy and E.~M.~Wright, \emph{An Introduction to the Theory of Numbers}, 6th ed., Oxford, 2008, Chapter~21, for $g(4)\ge19$ and $G(4)\ge16$.
\item \emph{The circle method.} G.~H.~Hardy and J.~E.~Littlewood, ``Some problems of `Partitio Numerorum' I: A new solution of Waring's problem'', \emph{Nachr.\ Ges.\ Wiss.\ Göttingen} (1920), 33--54. L.-K.~Hua, ``On Waring's problem'', \emph{Quart.\ J.\ Math.\ Oxford} 9 (1938), 199--202. R.~C.~Vaughan, \emph{The Hardy--Littlewood Method}, 2nd ed., Cambridge, 1997, Chapters~2--4: Theorem~4.6 there is Theorem~\ref{thm:singpos} for $s\ge4k$, and Lemmas~2.12--2.15 there are the local lemmas; our Lemma~\ref{lem:localsol} uses a cruder count, enough for $4k$. H.~Davenport, \emph{Analytic Methods for Diophantine Equations and Diophantine Inequalities}, 2nd ed., Cambridge, 2005.
\end{itemize}
